\chapter{Normalizing Flows}

\section{Change of Variables Formula}

Normalizing flows provide a class of generative models based on invertible transformations of probability distributions.

\medskip

\noindent
\textbf{Core idea.}  
Start with a simple base distribution $z \sim p(z)$ and transform it into a complex distribution $x = f_\theta(z)$.

\medskip

\noindent
\textbf{Change of variables formula.}  
If $f$ is invertible and differentiable, then
\[
p_\theta(x)
=
p(z) \left| \det \left( \frac{\partial z}{\partial x} \right) \right|,
\quad \text{where } z = f^{-1}(x).
\]

\medskip

\noindent
\textbf{Equivalent form.}
\[
\log p_\theta(x)
=
\log p(z)
+
\log \left| \det \left( \frac{\partial z}{\partial x} \right) \right|.
\]

\medskip

\noindent
\textbf{Interpretation.}
\begin{itemize}
    \item Density changes according to local volume distortion
    \item The determinant measures expansion or contraction
\end{itemize}

\medskip

\noindent
\textbf{Insight.}  
Normalizing flows transform simple distributions into complex ones while keeping track of density changes.

\section{Invertible Transformations}

The design of normalizing flows relies on constructing invertible transformations that are both expressive and computationally efficient.

\medskip

\noindent
\textbf{Requirements.}
\begin{itemize}
    \item $f$ must be invertible
    \item The Jacobian determinant must be tractable
\end{itemize}

\medskip

\noindent
\textbf{Composition of transformations.}
\[
x = f_L \circ f_{L-1} \circ \cdots \circ f_1(z).
\]

\medskip

\noindent
\textbf{Jacobian factorization.}
\[
\log \left| \det \frac{\partial z}{\partial x} \right|
=
\sum_{\ell=1}^L
\log \left| \det \frac{\partial h_{\ell-1}}{\partial h_\ell} \right|.
\]

\medskip

\noindent
\textbf{Types of transformations.}
\begin{itemize}
    \item Affine coupling layers
    \item Autoregressive transformations
    \item Volume-preserving flows
\end{itemize}

\medskip

\noindent
\textbf{Tradeoff.}
\begin{itemize}
    \item Expressivity vs computational efficiency
\end{itemize}

\medskip

\noindent
\textbf{Insight.}  
Careful design of invertible mappings enables tractable and expressive generative models.

\section{Exact Likelihood Computation}

A key advantage of normalizing flows is the ability to compute exact likelihoods.

\medskip

\noindent
\textbf{Likelihood.}
\[
\log p_\theta(x)
=
\log p(f^{-1}(x))
+
\log \left| \det \left( \frac{\partial f^{-1}}{\partial x} \right) \right|.
\]

\medskip

\noindent
\textbf{Comparison with other models.}
\begin{itemize}
    \item VAEs: approximate likelihood
    \item GANs: implicit likelihood
    \item Flows: exact likelihood
\end{itemize}

\medskip

\noindent
\textbf{Sampling.}
\begin{itemize}
    \item Sample $z \sim p(z)$
    \item Compute $x = f_\theta(z)$
\end{itemize}

\medskip

\noindent
\textbf{Inference.}
\begin{itemize}
    \item Compute $z = f_\theta^{-1}(x)$
    \item Evaluate density using change of variables
\end{itemize}

\medskip

\noindent
\textbf{Efficiency considerations.}
\begin{itemize}
    \item Jacobian determinant must be efficiently computable
    \item Often achieved via triangular or structured matrices
\end{itemize}

\medskip

\noindent
\textbf{Insight.}  
Normalizing flows provide a rare combination of exact likelihood and efficient sampling.

\section{Model Design}

Designing effective normalizing flows involves balancing expressivity and tractability.

\medskip

\noindent
\textbf{Coupling layers.}
\begin{itemize}
    \item Partition variables into two parts
    \item Transform one part conditioned on the other
\end{itemize}

\medskip

\noindent
\textbf{Autoregressive flows.}
\begin{itemize}
    \item Model dependencies sequentially
    \item Triangular Jacobian simplifies determinant computation
\end{itemize}

\medskip

\noindent
\textbf{Stacking transformations.}
\begin{itemize}
    \item Multiple layers increase expressivity
    \item Each layer contributes to transformation complexity
\end{itemize}

\medskip

\noindent
\textbf{Invertible neural networks.}
\begin{itemize}
    \item Architectures designed to ensure invertibility
\end{itemize}

\medskip

\noindent
\textbf{Limitations.}
\begin{itemize}
    \item Architectural constraints for invertibility
    \item Computational cost for high-dimensional data
\end{itemize}

\medskip

\noindent
\textbf{Insight.}  
Model design in flows requires balancing flexibility with computational constraints.

\section{A Unifying Perspective}

Normalizing flows combine several key ideas:

\begin{itemize}
    \item \textbf{Transformations:} mapping simple distributions to complex ones
    \item \textbf{Invertibility:} enabling exact density computation
    \item \textbf{Composition:} building expressive models through layers
\end{itemize}

\medskip

\noindent
\textbf{Core components.}
\begin{itemize}
    \item Base distribution $p(z)$
    \item Invertible mapping $f_\theta$
    \item Jacobian determinant
\end{itemize}

\medskip

\noindent
\textbf{Key insight.}  
By carefully designing invertible transformations, we can build expressive generative models with exact likelihoods.

\section{Outlook}

This chapter introduced normalizing flows:
\begin{itemize}
    \item change of variables for density transformation,
    \item invertible transformations,
    \item exact likelihood computation,
    \item practical model design.
\end{itemize}

\medskip

In the next chapter, we study diffusion models, which approach generative modeling through stochastic processes and iterative refinement.

\medskip

\noindent
\textbf{Guiding principle.}  
Structured transformations of probability distributions provide a powerful framework for generative modeling.